\section{Estimates for class M2: the nonshort cases}\label{app:nonshort-estimates}

This appendix proves the scalar inequalities behind \cref{sec:mixed}, in the notation of
\cref{tab:m2-notation}. Every step concerns explicit functions of at most five real variables.
The analytic ingredients are a few rational bounds for logarithms; everything else is polynomial,
and every polynomial sign is reduced, one variable at a time, to univariate inequalities on
intervals. Several of the polynomials are long; they are determined by the displayed definitions,
and we state the identities and bounds on which each sign rests. The arguments follow
\cite{Notes} and, where indicated, the shorter arguments of \cite{LighterProofs}.

\subsection{Secants and floors}\label{app:m2-secants}

For $B>0$ the function $g_B(x)=1-Bx/(B+x)^2$ decreases on $(0,B)$, increases on $(B,\infty)$, and
has minimum $g_B(B)=3/4$; also $g_B(x)=U(B/(B+x))=U(x/(B+x))$. The turn secant
$\operatorname{Sec}(A,B,P)=(P-A)^{-1}\int_A^Pg_B$ has the antiderivative
$s-B\log(B+s)-B^2/(B+s)$.

\begin{lemma}[Secants]\label{lem:m2-secants}
\begin{enumerate}[label=(\alph*)]
\item If $0<a\le c<d$, then $\operatorname{Sec}(a,c,d)<\operatorname{Sec}(a,d,c)$.
\item If $0<c\le d$, then $z\mapsto\operatorname{Sec}(z,d,c)$ strictly decreases on $(0,c]$.
\item For $P,C>0$ and $q>1$, $F_{P/(P+C)}(q)=\operatorname{Sec}((P+C)/q-P,\,P,\,C)$. In particular
$F_{d/W}(q)=\operatorname{Sec}(B_0,d,c)$ with $B_0=(4-dr)/(d+r)$, the third margin at the switch
$d\to r$.
\item If $ad<c^2$, then $\operatorname{Sec}(a,c,d)<\frac12(g_c(a)+g_c(d))$.
\item If $B/(A+B)\ge3/5$ and $A<P$, then $\operatorname{Sec}(A,B,P)>94/125$.
\end{enumerate}
\end{lemma}

\begin{proof}
(a) $\operatorname{Sec}(a,d,c)=\int_0^1g_1(u)\,d\theta$ and $\operatorname{Sec}(a,c,d)=\int_0^1g_1(v)\,d\theta$
with $u=(a+\theta(c-a))/d$ and $v=(a+\theta(d-a))/c$. Here $v>u$, $uv<1$, and
$g_1(u)-g_1(v)=(v-u)(1-uv)/((1+u)^2(1+v)^2)>0$.
(b) The derivative is $(\operatorname{Sec}(z,d,c)-g_d(z))/(c-z)$, and the average of the decreasing
function $g_d$ over $[z,c]$ is below its value at $z$.
(c) Substitute $(P+C)/q-P$ in the closed form of the secant; $W/q-d=Q/Z-d=B_0$.
(d) For $a\le x\le d$,
\begin{multline*}
\frac{(d-x)g_c(a)+(x-a)g_c(d)}{d-a}-g_c(x)\\
=\frac{c(x-a)(d-x)\,[c^2(a+d+2c)+(c^2-ad)x]}{(a+c)^2(d+c)^2(x+c)^2}\ge0;
\end{multline*}
integrate over $[a,d]$.
(e) With $\eta=B/(A+B)$ and $q=(P+B)/(A+B)>1$, the secant is $1+\eta^2/q-\eta\log q/(q-1)$.
Bound $\log q/(q-1)$ above by a rational function of $q$ as in \cref{lem:m2-logs}(b); the
resulting rational lower bound exceeds $94/125$ for $\eta\ge3/5$, by a discriminant in $q$.
\end{proof}

\subsection{Rational bounds for the logarithms}\label{app:m2-logs}

\begin{lemma}[Logarithms]\label{lem:m2-logs}
Let $x>0$, $q>1$, $c=1-t$, $W=1+x$, $X=x+2$, $p=t+x$, $g=1-cp/W^2$, and let
$v=1+c^2/W-c\log W/x$ be the normalized incoming secant.
\begin{enumerate}[label=(\alph*)]
\item $\log(1+x)>2x/X+2x^3/(3X^3)$; hence $v<\bar v=1+c^2/W-(2c/X)(1+x^2/(3X^2))$ when $t<1$.
\item $\log q<L_+(q)=(q-1)(q^2+10q+1)/(6q(q+1))$; hence
$F_{p/W}(q)>F_\ell=1+p^2q/W^2-(p/W)(q^2+10q+1)/(6(q+1))$.
\item $\log(1+x)>x(43x^3+252x^2+390x+180)/(12W^2(x^2+10x+15))$; hence
$v<v_G=g+(cx/W^2)(T_G(x)-t)$ with $T_G=(90-19x^2)/(12(x^2+10x+15))$.
\item $\log(1+x)>3x(x+2)/(x^2+6x+6)$; hence $v<V=1+c^2/W-3c(x+2)/(x^2+6x+6)$.
\item $\bar v-g=((1-t)x/W^2)(T(x)-t)$ with $T(x)=(-2x^3-x^2+10x+12)/(3(x+2)^3)$; so $v\ge g$ forces
$t<T(x)$. The function $T$ decreases, and $T(1)=19/81<1/4$.
\end{enumerate}
\end{lemma}

\begin{proof}
(a) Integrate $1/(1-z^2)>1+z^2$ from $0$ to $x/X$. In (b), (c), (d) the difference vanishes at the
left end and has derivative $(q-1)^4/(6q^2(q+1)^2)$, $x^6/(W^3(x^2+10x+15)^2)$ and
$x^4/(W(x^2+6x+6)^2)$ respectively. (e) Direct subtraction; $T'(x)=-(11x^2+24x+16)/(3(x+2)^4)$.
\end{proof}

\subsection{The residual}\label{app:m2-residual}

Put $c=1-t$, $p=t+x$, $W=1+x$, $Y=x(x+2)$, $B=k+2tx+x^2$, $Z=qB/W$ (the normalized $d+r$), and
\[
H_0=\frac kt\,(v-1+t^2)+t(2-t),\qquad w_U=\frac{k^2-kt+t^2+H_0Y+vY^2}{B^2},
\]
the central diagonal and the transported outgoing norm with the incident norm at the boundary
ceiling $U(t/k)$. The root function at the ceiling, divided by $T$, is
\[
R=c+Z-\frac{p^2}{W}-\frac{p^2}{Z}+\frac{p^3}{B}-p\log q-Wv-Zw_U+\frac{p(H_0+2vY)}{B},
\]
and $N=tBqW\,R$. Clearing denominators cancels every negative power, and $N$ is a polynomial,
affine in $k$, $v$ and $L=\log q$:
\begin{itemize}
\item the coefficient of $v$ is $N_v=-xq\,I$ with
$I=k+tx+tW(k-2t)+(q-1)\bigl(kX+txX^2\bigr)$, nonpositive when $k\ge2t$;
\item the coefficient of $L$ is $N_L=-tpqWB\le0$;
\item the coefficient of $k$ is
\begin{multline*}
A=q\bigl\{x[(1-v)(1+tW)-t(1-t^2)]\\
+(q-1)[t^2(1+2x-x^2)+2tx^2+(1-v)Y]-tpW L\bigr\}.
\end{multline*}
\end{itemize}
In $A$ the coefficient of $v$ is $-qx[(q-1)(x+2)+1+tW]<0$ and that of $L$ is $-qtpW<0$. Inserting
$v\le U(t)=1-t+t^2$ and $L<q-1$ and collecting gives
$A>q\,t\,x(1-t)[(q-1)(2x+1)+tx]>0$, which is \cref{lem:m2-shape}.

\begin{lemma}[Rational relaxation]\label{lem:m2-relax}
If $k\ge2t$, $q>1$ and $F_{p/W}(q)\le v$, then
\[
N(v,\log q)\ \ge\ N(\bar v,L_+)-(q-1)(\bar v-F_\ell),
\]
and $P_k:=6(q+1)W^2X^3\{N(\bar v,L_+)-(q-1)(\bar v-F_\ell)\}$ is a polynomial in $t$, $x$, $k$ and
$y=q-1$, cubic in $y$.
\end{lemma}

\begin{proof}
$N_v\le0$, $N_L\le0$, $v<\bar v$ and $\log q<L_+$, while $F_\ell<F_{p/W}(q)\le v<\bar v$ makes the
subtracted term nonnegative.
\end{proof}

\subsection{The necessary bounds}\label{app:m2-bounds}

We prove \cref{lem:m2-necessary}(iii)--(vii).

\runin{The scale.} At the centre, \eqref{eq:m2-centre} holds with $u_{\mathrm i}=u$,
$u_{\mathrm o}=v$, $r=a/S$ and $t=a/T$. Lowering both norms to a floor $\mu$ and raising both ratios
to a cap $\zeta$ gives $\gamma^2\ge(4\mu^2-B(\rho)^2)/(16(1-\rho)^2)$ with
$B(\rho)=1-\rho+\log\rho+2\psi_\mu(\zeta)$, $\psi_\mu(\zeta)=\zeta-\log\zeta-(1-\mu)/\zeta$. The
cleared comparison $\varphi(\rho)=4\mu^2-B(\rho)^2-16\,\Gamma_{\max}(1-\rho)^2$, with
$\Gamma_{\max}=9193/50000$, has derivative $2(1-\rho)(16\Gamma_{\max}-B/\rho)$ and increases, so its
positivity at the left end of an interval of $\rho$ excludes that interval by
\cref{lem:m2-casimir}(b). With the floor $3/4$ and the cap $1/2$ this gives $a^2<2/3$, that is
$\rho<1/7$. When $t\le2/5$, both ratios are at most $2/5$ and both turns have middle-margin ratio
$B/(A+B)\ge3/5$, so \cref{lem:m2-secants}(e) gives $\mu=94/125$; the interval $[4/29,3/20)$ is then
excluded, and $a<4/5$. In normalized form $a^2(k-t^2)=4t^2$, so $a^2<2/3$ is $k>7t^2$, and $a<4/5$
is $k>29t^2/4$.

\begin{lemma}[The ratio $t<2/5$]\label{lem:m2-high-t}
Every window of \cref{prop:m2-window} has $t<2/5$.
\end{lemma}

\begin{proof}
Suppose $t\ge2/5$. Then $k>7t^2\ge28/25$. By \cref{lem:m2-logs}(e), $t<T(x)$, and the numerator
$12-22x-41x^2-16x^3$ of $T(x)-2/5$ is negative from $x=1/3$ on, so $x<1/3$. The centre gap (below)
gives $x^2+2x>7t^2-2t\ge8/25$, so $x>1/7$. Lower $k$ to $7t^2$ (\cref{lem:m2-shape}) and use
\cref{lem:m2-relax}: $P_{7t^2}=C_0+C_1y+C_2y^2+C_3y^3$ with $C_j=C_j(t,x)$. Each $C_j$ increases in
$t$ on $2/5\le t\le1/2$, $0<x\le1/3$. Indeed $\partial_tC_j$ is a polynomial in $x$ whose
coefficients are cubic in $t$; the bounds $\frac45t-\frac4{25}\le t^2\le\frac9{10}t-\frac15$ and
$\frac{12}{25}t-\frac{16}{125}\le t^3\le\frac{61}{100}t-\frac9{50}$ (with remainders $(t-\frac25)^2$,
$(t-\frac25)(\frac12-t)$, $(t-\frac25)^2(t+\frac45)$, $(t-\frac25)(\frac12-t)(t+\frac9{10})$) make every
coefficient affine in $t$, and the resulting polynomials in $x$ are positive. At $t=2/5$,
\[
C_0=-\tfrac{72}{625}x^2(x+1)(15x^4+68x^3+94x^2+20x-48),
\]
\[
C_1=\tfrac{12x}{625}(84+2086x-395x^2-8726x^3-11540x^4-6493x^5-1647x^6-140x^7),
\]
\begin{multline*}
C_2=-\tfrac6{625}(144-1444x-3342x^2+25127x^3+75813x^4\\
+88565x^5+54334x^6+18396x^7+3270x^8+250x^9),
\end{multline*}
\begin{multline*}
C_3=\tfrac{x+2}{625}(2680+17612x+33046x^2+3121x^3-42903x^4\\
-45902x^5-20182x^6-3790x^7-250x^8).
\end{multline*}
Here $C_0>0$ for $x\le1/3$. The function $C_1$ has at most one interior critical point, a maximum,
and $C_2$ at most one, a minimum, because $C_1'/x$ and $C_2'/x$ are monotone; and $C_3$ increases.
Evaluation at the endpoints therefore gives, on the rows $[\frac17,\frac15]$, $[\frac15,\frac3{10}]$,
$[\frac3{10},\frac13]$, the bounds $C_1>b$, $C_2>-M$, $C_3>c$ with
$(b,M,c)=(\frac9{10},1,20)$, $(\frac32,10,26)$, $(\frac{77}{50},\frac{771}{50},\frac{194}5)$, and
$4bc>M^2$ in each row. Hence $C_1y+C_2y^2+C_3y^3\ge y(b-My+cy^2)\ge0$, so $P_{7t^2}>0$ and $N>0$,
against $N<0$.
\end{proof}

\runin{The centre gap $d^2+2cd>4$.} With $u$ at its ceiling, the diagonal and the transported norm
increase in $b$: $\partial_bh_U=K_a(b-a)/(16S^2)$ and $\partial_bw_U=K_a(b-a)(aK_a+SY)/(S^3Q^2)$. So
the actual outgoing norm is below its value $w_U(b_*)$. The incoming secant is below the chord
value $V_{\mathrm{ch}}=\frac12(g_c(a)+g_c(d))$ (\cref{lem:m2-secants}(d)), and $w_U$ increases in
$v$. If $e=4-d^2-2cd\ge0$, a cleared identity with five factored coefficients puts $w_U(b_*)$ at
$v=V_{\mathrm{ch}}$ strictly below $g_d(r_0)\le g_d(r)\le w$, with $r_0=Q/W-d$, which is impossible.

\runin{The duration $3d<7c$.} Since $v=\operatorname{Sec}(a,c,d)\le g_c(a)$ and $g_c\ge g_c(a)$ on
$[0,a]$, extending the interval to $[0,d]$ raises the average:
$\operatorname{Sec}(0,c,d)>v\ge g_c(d)$. With $h=d/c$ this reads
$1+1/(1+h)-\log(1+h)/h>(1+h+h^2)/(1+h)^2$, which forces $h<7/3$; at $h=7/3$ the two sides are
$1.3-\frac37\log\frac{10}3<0.785$ and $0.79$.

\runin{The reverse order $d\le b_*$.} Suppose $d>b_*\ge c$. By \cref{lem:m2-secants}(a),(c),
$\operatorname{Sec}(a,d,c)>\operatorname{Sec}(a,c,d)\ge\operatorname{Sec}(B_0,d,c)$, and
\cref{lem:m2-secants}(b) forces $B_0>a$, that is $(a+d)(a+r)<K_a$; with $d>b_*$ this gives $r<c<d$.
The reverse successor gives $w\le\operatorname{Sec}(a,r,d)$, and the outgoing floor is
$\operatorname{Sec}(A_0,d,r)$ with $A_0=(4-cd)/(c+d)$. The same two steps force $A_0>a$, but
$d>b_*$ is equivalent to $A_0<a$. In normalized form $b_*/T=k-t$ and $d/T=t+x$, so $d\le b_*$ is
$k\ge2t+x$.

\subsection{The reverse successor}\label{app:m2-reverse-signs}

\begin{lemma}[Unit shape]\label{lem:m2-unit-shape}
If $0<t\le13/35$, $0<x\le1-2t$, $2t+2x+x^2>1$, $q>1$, $v=\operatorname{Sec}(t,1-t,t+x)$ and
$F_{p/W}(q)\le v$, then $N(t,x,1,q,v,\log q)>0$.
\end{lemma}

\begin{proof}
By \cref{lem:m2-relax} it suffices that $P_1=C_0+C_1y+C_2y^2+C_3y^3>0$ for $y\ge0$. This holds when
\[
C_0>0,\qquad J_1=20xC_1+3C_0>0,\qquad J_2=53C_0+20x^2C_2>0,
\]
\[
J_3=24x^3C_3+12x^2C_2-C_0>0,\qquad J_4=40x^3C_3+30x^2C_2+7C_0>0,
\]
\[
J_*=24x^3C_3-C_0>0.
\]
Indeed, with $z=y/x$ and $B_1=-x^2C_2/C_0$, these give $xC_1/C_0\ge-\frac3{20}$, $B_1\le\frac{53}{20}$
and $x^3C_3/C_0\ge\max\{\frac1{24},\frac{B_1}2+\frac1{24},\frac{3B_1}4-\frac7{40}\}$. The resulting lower
bound for $P_1/C_0$ is affine in $B_1$ on each side of $B_1=13/15$, so it suffices to take
$B_1\in\{0,\frac{13}{15},\frac{53}{20}\}$, that is, to check the cubics $1-\frac3{20}z+\frac1{24}z^3$,
\[
1-\tfrac3{20}z-\tfrac{13}{15}z^2+\tfrac{19}{40}z^3=\tfrac{52}{135}+\tfrac{13}{180}(z-\tfrac43)+\tfrac25(z-\tfrac43)^2+\tfrac{19}{40}z(z-\tfrac43)^2,
\]
\[
1-\tfrac3{20}z-\tfrac{53}{20}z^2+\tfrac{29}{16}z^3=\tfrac1{80}-\tfrac1{80}(z-1)+\tfrac{39}{40}(z-1)^2+\tfrac{29}{16}z(z-1)^2,
\]
each bounded below by a positive constant on $z\ge0$.

For the six signs put $s=1-2t-x\ge0$. The domain gives $\frac3{25}<x<1$ and
$\max\{0,\frac9{35}-x\}\le s\le\min\{1-x,x+x^2\}$, and $s=\frac9{35}-x$ is the edge $t=13/35$.

First, $C_0=12W^2X^3N(\bar v,L_+)|_{q=1,k=1}$. At $q=1$ the residual is affine in $k$ with
$\partial_k=x\{(1-\bar v)(1+tW)-t(1-t^2)\}>t^2(1-t)x^2>0$, and $2t+x\le1$, so
\[
C_0\ \ge\ 12W^2X^3\,\frac{(1-t)x^3}{3WX^3}\,P_0(t,x),
\]
\[
P_0=t(3t-1)x^3+(18t^2-15t+5)x^2+(36t^2-40t+14)x+(24t^2-32t+12),
\]
where the right side is the value at $k=2t+x$. From $t(3t-1)\ge-\frac1{12}$, the completed squares
$18(t-\frac5{12})^2+\frac{15}8$, $36(t-\frac59)^2+\frac{26}9$, $24(t-\frac23)^2+\frac43$, and $x<1$, one gets
$P_0>\frac{43}{24}x^2+\frac{26}9x+\frac43>0$.

Next write $J_i/x^2=\sum_{m=0}^3A_{im}(x)s^m$ for $i=1,\dots,4$. The functions $J_1/x^2$ and $J_2/x^2$
increase in $s$: their $s$-derivatives are bounded below, using $s\le x+x^2$, by polynomials in $x$
with positive coefficients. So they are smallest at the lower end of $s$, where they are positive:
on the edge $t=13/35$ by one-variable inequalities on $[\frac3{25},\frac9{35}]$, and at $s=0$ because
$A_{i0}/(xW^2)$ increases for $x\ge\frac9{35}$ and is positive at $\frac9{35}$. The function $J_3/x^2$ is
concave in $s$ ($A_{32}+3A_{33}s<0$ on the domain), so it suffices to check the two ends of the
$s$-interval, again one-variable inequalities. The function $J_4/x^2$ increases in $s$ for
$x\ge\frac18$: its $s$-derivative, written in the Bernstein basis of $[0,x+x^2]$, has three
coefficients that are polynomials in $x$ with positive non-constant coefficients and positive at
$x=\frac18$. For $\frac3{25}<x<\frac18$, one has $A_{40}\ge\frac{7059}{800}$, $A_{41},A_{43}>0$,
$A_{42}>-217$ and $s\le\frac9{64}$, so $J_4/x^2>4$. Finally $J_*$, a cubic in $s$ written in the
Bernstein basis of $[0,S]$ with $S=x+x^2$ for $x\le\frac25$ and $S=1-x$ for $x\ge\frac25$, has four
coefficients bounded below, after dropping positive terms, by $258xW^2$, $675xW^2$, $570xW^2$,
$94xW^2$ in the first case and by $258xW^2$, $100W$, $46$, $94$ in the second.
\end{proof}

\begin{lemma}[Order face]\label{lem:m2-order-face}
If $0<t<2/5$, $x>1/4$, $2t+x\ge1$, $x<\frac73-\frac{10}3t$, $v=\operatorname{Sec}(t,1-t,t+x)$ and
$1<q\le(1+x)/(2t+x)$, then $N(t,x,2t+x,q,v,\log q)>0$.
\end{lemma}

\begin{proof}
Put $k=2t+x$. Then $B=Wk$, $Z=qk$, $r=Z-p$, and with $t,x,k,v$ fixed,
$R'(q)=k(g_p(r)-w_U)$ and $R''(q)=k^2p(r-p)/(r+p)^3\le0$, since $t\le r\le c\le p$ on the interval.
So $R$ is concave, and it suffices to check $q=1$ and $q=W/k$.

At $q=1$, $N=x\{(k+tY)(1-v)-2t(1-t)(t+x)\}$ with $k+tY=tW^2+t+x$, and \cref{lem:m2-logs}(a) gives
$N>\frac{(1-t)x^3}{3WX^3}P_0(t,x)$ with $P_0$ as in \cref{lem:m2-unit-shape}, positive (for $t<\frac13$
use $t(1-3t)\le\frac1{12}$ and $x<\frac52$).

At $q=W/k$ replace $v,\log q$ by $\bar v,L_+$, which lowers $R$. With $h=1-2t$, $s=k-1$ and $x=h+s$
the domain is $\frac15<h<1$, $0\le s<\frac43$, $h+s>\frac14$, and
\[
\bar R(W/k)=\frac{P(h,s)}{24tk^2W^2(W+k)X^3},\qquad P=\sum_{j=0}^7C_j(h)s^j,
\]
\[
C_0=h^3(1+h)^2F(h),\qquad F=-16+24h+120h^2+140h^3+78h^4+23h^5+3h^6,
\]
\[
C_1=h^3(1+h)(8+396h+1020h^2+1080h^3+597h^4+179h^5+24h^6),
\]
and $C_2,\dots,C_7$ are polynomials in $h$ whose only negative monomials are $-70h^2$ in $C_5$,
$-108h^2$ in $C_6$, and $-4h$, $-20h^2$ in $C_7$. These are absorbed: $184hs^5-70h^2s^5\ge0$, and
$316hs^4-108h^2s^6-(4h+20h^2)s^7\ge hs^4(316-108\cdot\frac{16}9-24\cdot\frac{64}{27})>0$, the term
$316hs^4$ being part of $C_4s^4$. So $P\ge C_0+C_1s$. The polynomial $F$ increases and
$F(\frac14)=\frac{63}{4096}>0$, which settles $h\ge\frac14$. For $\frac15\le h\le\frac14$ one has
$s>\frac14-h$ and $C_0+C_1(\frac14-h)=\frac{h^3(1+h)}4Q(h)$ with
$Q=-56+396h+12h^2-1960h^3-2851h^4-1805h^5-588h^6-84h^7$, which is concave on $[\frac15,\frac14]$
($Q''<24-\frac{11760}5<0$), with $Q(\frac15)=\frac{220476}{78125}$ and $Q(\frac14)=\frac{315}{4096}$.
\end{proof}

\begin{lemma}[Scale strip]\label{lem:m2-strip}
If $\frac{13}{35}\le t\le\frac25$, $0<x\le\frac9{35}$, $\frac{29}4t^2<2t+2x+x^2$, $q>1$,
$v=\operatorname{Sec}(t,1-t,t+x)$ and $F_{p/W}(q)\le v$, then $N(t,x,\frac{29}4t^2,q,v,\log q)>0$.
\end{lemma}

\begin{proof}
Write $P_{29t^2/4}=\sum_jC_j(t,x)y^j$. Each $C_j$ increases in $t$ on the strip:
$\partial_tC_0\ge44x^2$, $\partial_tC_1\ge99x$, $\partial_tC_2\ge48-\frac{2301}{81}$ and
$\partial_tC_3\ge64-\frac{167}{729}$. Each $\partial_tC_j$ is a polynomial in $x$ whose coefficients
are cubics $\sum_ia_it^i$, and on $[\frac{13}{35},\frac{14}{35}]$ such a cubic is at least
$\sum_{a_i\ge0}a_i(\frac{13}{35})^i+\sum_{a_i<0}a_i(\frac{14}{35})^i$; one coefficient of $\partial_tC_0$
needs the additional remark that it decreases in $t$. So
$P_{29t^2/4}(t,x,y)\ge P_{29t_0^2/4}(t_0,x,y)$ with $t_0=\frac{13}{35}$, for every $y\ge0$.
Since $\frac{29}4t^2-2t$ increases, the strip condition gives
$2x+x^2>\frac{29}4t_0^2-2t_0>\frac9{35}$, so $(t_0,x)$ satisfies the shape inequalities of
\cref{lem:m2-unit-shape}, whose polynomial part gives $P_1(t_0,x,y)>0$ (no floor is used there). As
$P_k$ increases in $k$ at fixed $(t_0,x,q)$ (the coefficient of $k$ is bounded below as in
\cref{lem:m2-shape}, now with $\bar v<U(t_0)$ and $L_+<q-1$) and $\frac{29}4t_0^2=\frac{4901}{4900}>1$,
also $P_{29t_0^2/4}(t_0,x,y)>0$. Hence $P_{29t^2/4}>0$, and \cref{lem:m2-relax} at the original
$(t,x,q)$, where the floor holds, gives $N>0$.
\end{proof}

\begin{lemma}[Low-scale unit face]\label{lem:m2-low-unit}
If $0<t\le\frac14$, $\frac12<x\le\frac73$, $x\ge1-2t$, $t\le T(x)$, $q>1$,
$v=\operatorname{Sec}(t,1-t,t+x)$ and $F_{p/W}(q)\le v$, then $N(t,x,1,q,v,\log q)>0$.
\end{lemma}

\begin{proof}
With $I_0=1+tx+tW(1-2t)$, $I_1=X+txX^2$ and $\mathcal A=xq(I_0+I_1y)=-N_v$, one has
$N>N(\bar v,L_+)-\mathcal A(\bar v-F_\ell)$, and
$Q=6(q+1)W^2\{N(\bar v,L_+)-\mathcal A(\bar v-F_\ell)\}$ is a quartic $P_0+P_1y+\dots+P_4y^4$. It is
positive for $y\ge0$ when $P_0>0$, $P_1+2P_0>0$, $P_2+48P_0>0$, $P_3+4P_2>0$, $P_3-P_0>0$ and
$P_4-P_0>0$; the extreme cases reduce to
$1-2y-\frac14y^2+y^3+y^4=\frac4{81}+\frac5{27}(y-\frac23)+(\frac{29}{12}+\frac73y+y^2)(y-\frac23)^2$ and
$1-2y-48y^2+192y^3=\frac29-2(y-\frac16)+(16+192y)(y-\frac16)^2$, both positive. The first four signs
factor as $4(1-t)x^2F_0/X^3$, $2(1-t)xF_1/X^3$, $F_2/X^3$, $F_3/X^3$ with polynomials $F_i(t,x)$, and
$-\partial_tF_i$ is smallest at $t=\frac14$ and positive there (in the variables $1-4t$ and $2x-1$ its
coefficients are nonnegative after absorbing a few negative terms), so
$F_i(t,x)\ge F_i(\frac14,x)+(\frac14-t)(-\partial_tF_i(\frac14,x))$. For $\frac12\le x\le1$ the values
$F_i(\frac14,x)$ are positive one-variable polynomials. For $1\le x\le\frac73$, where
$t<T(x)<\frac14$, the tangent value at $t=T(x)$ is a positive multiple of an explicit polynomial
$h_i(x)$, positive on $[1,\frac73]$ after charging its few negative high-degree terms to positive
lower ones. The comparisons $P_3>P_0$ and $P_4>P_0$ are proved in the same way.
\end{proof}

\subsection{The forward successor}\label{app:m2-forward-signs}

In this subsection write $o=2t+x$, $j=o-1$, $h=1-2t$, $y=q-1$, $e=q/W$, $J=t+Y$, $D=eJ-p$ and
$E=e^2B+1-2pe$. The next norm at the ceiling is $\mathcal T_U=(v-e^2(H_0+2vY)+e^4B^2w_U)/E^2$, and
the combined residual \eqref{eq:m2-comb} has $C_*=tW^2DE^2/e$; it is a concave quadratic in $k$
with leading coefficient $-tW^2De^3/4$.

\begin{lemma}[Low ratio]\label{lem:m2-low-ratio}
Let $\frac14\le t\le\frac13$, $1-2t\le x\le1$, $1\le k\le o$, $1\le q\le W/o$, and let
$\widehat G=\mathcal H(\hat v,L_+)$, with $\alpha=6$, $M=0$, $\hat v=v_G+\rho(v_G-F_\ell)$ and
$\rho=\frac34(q+o-1)$. Then $\widehat G>0$. Consequently every window with a forward successor has
$t>1/3$.
\end{lemma}

\begin{proof}
At fixed $(t,x,q)$ the coefficients of $v$ and $\log q$ in $\mathcal H$ are negative:
$\mathcal H_v<-\frac{17}{12}xq(I_0+I_1y)$, from $Je^2\le\frac{10}3$ and $WD/x\le\frac94$, and
$\mathcal H_{\log q}=(7-Je^2)N_L<0$. So a window with a forward successor and $t\le\frac13$, which
lies in the domain (with $q<W/(1+2j)\le W/o$), has $\widehat G\le\mathcal H(v,\log q)<0$.

The three-term secant majorant $v_3=1+c^2/W-(2c/X)(1+x^2/(3X^2)+x^4/(5X^4))$ satisfies
$v_3-v_G=cx^6(31x^2+82x+46)/(60W^2X^5(x^2+10x+15))>0$, and $v_3+\rho(v_3-F_\ell)<U(t)$ (using
$F_\ell-g=\eta(\eta-\frac12)y+\eta y^2/(3(y+2))\ge0$ with $\eta=p/W\ge\frac12$, and $\rho<2$); so
$\hat v<U(t)$. Also $\mathcal H_{kv}<0$ and $\mathcal H_{k,\log q}<0$, so
$\widehat G_k>\mathcal H_k(U(t),y)$, which decreases in $k$. At $k=o$,
\[
\mathcal H_k(U(t),y)\big|_{k=o}=\frac{tq}{2W^2}(d_0+d_1y+d_2y^2+d_3y^3),
\]
with $d_0=xc\{1-2t+x(-2t^2+14t-1)+24tx^2+12tx^3\}>0$, $d_3=xWJ>0$, $d_2<0$ (every coefficient
in $x$ is negative for $\frac14\le t\le\frac13$), and $2d_1+d_2>0$; as $y\le\frac12$,
$d_1+d_2y\ge d_1+\frac12d_2>0$. Hence $\widehat G_k>0$, and it suffices to show $\widehat G>0$ at $k=1$.

For $k=1$ we follow \cite[\texttt{NONSHORT\_\allowbreak LOW\_\allowbreak RATIO\_\allowbreak LINK}]{LighterProofs}.
With $t=\frac{1-h}2$, $x=h+j$ and $q=1+y$ on $\frac13\le h\le\frac12$, $j\ge0$, $0\le y\le h$, a region
that contains the domain, the polynomial $\Pi=192q(q+1)W^3(x^2+10x+15)\,\widehat G(t,x,q,1)$ equals
$\sum_i\tilde C_i(h,y)j^i$. Its part of degree at most two in $j$ is
$h^3w^3(\mathcal A_0+\beta\mathcal A_1+w\beta^2\mathcal A_2)$ with $w=1+h$ and $\beta=j/(hw)$; each
$\mathcal A_i$ increases in $h$, its $h$-derivative being a polynomial with nonnegative
coefficients after factoring, and at $h=\frac13$ the quadratic in $\beta$ exceeds $4$ by three square
identities. The coefficients $\tilde C_3,\dots,\tilde C_{10}$ are nonnegative from their factored
forms. Hence $\Pi>4h^3w^3+700j^3+3900j^4>0$.
\end{proof}

\begin{lemma}[Large gap]\label{lem:m2-large-gap}
Every window with a forward successor has $j<1/30$.
\end{lemma}

\begin{proof}
Suppose $j\ge\frac1{30}$. By \cref{lem:m2-low-ratio}, $t>\frac13$; \cref{lem:m2-logs}(c) with
$v\ge g$ gives $t<T_G(x)$, hence $x<\frac35$ and $7x<12h$, that is, $j<\frac57h$.

For $\frac13\le t\le\frac38$ we follow
\cite[\texttt{NONSHORT\_\allowbreak LOW\_\allowbreak T\_\allowbreak CALLER}]{LighterProofs}. Take
$\alpha=\frac72-10y\ge\frac{11}{16}$ and $M=8j+8y-\frac16\ge\frac1{10}$. The $v$-coefficient of $\mathcal H$
is negative, so $v\mapsto V_L$ (a rational majorant of $v$), $\log q\mapsto L_+$ and $F\mapsto F_\ell$
keep $\mathcal H<0$. The resulting residual increases in $k$ on $[1,o]$: it is quadratic in $k$,
its shape derivative at the midpoint of two shapes is affine in $v$, $\log q$ and $k$ with
nonpositive slopes, and at the extreme values it equals $q^2S/(12W)$ with a quartic $S$ in $y$ that
is concave in $y$ and positive at both ends. At $k=1$ a positive multiple $S$ of the residual, a
polynomial in $(h,j,y)$, increases in $h$ at fixed $(j,y)$: a cubic Taylor bound from
$\partial_h^4S\ge500000$ and two derivative signs at $h=\frac14$ give $\partial_hS>0$. So it suffices to
check the three faces $h=\frac14$, $j=\frac57h$ and the largest $y$, where $S>0$ by termwise monomial
bounds after a shift of the variable.

For $\frac38\le t<\frac25$ we follow
\cite[\texttt{NONSHORT\_\allowbreak HIGH\_\allowbreak RATIO\_\allowbreak SHAPE},
\texttt{NONSHORT\_\allowbreak HIGH\_\allowbreak RATIO\_\allowbreak LOWER\_\allowbreak REDUCTION}]{LighterProofs}.
Take $\alpha=4$, $M=1$ and $v\mapsto v_G$. The residual increases in $k$ on $(\frac{29}4t^2,o)$, by the
same midpoint argument with endpoint value $2tq^2(d_0'+d_1'y+d_2'y^2+d_3'y^3)$, a positive cubic.
At $k=\frac{29}4t^2$ a positive multiple $S$ satisfies $-\partial_h^2S>0$, and $-\partial_hS>0$ at
$h=\frac15$, which reduces the sign to $h=\frac14$; there $S$ is concave in $y$, and along the remaining
line a ninth-derivative bound and exact values at the ends give $S\ge\frac{13}4$. In both cases
$\mathcal H>0$, a contradiction.
\end{proof}

\begin{lemma}[The ratio $\chi$]\label{lem:m2-chi}
If $\frac13\le t\le\frac{19}{50}$, $0\le j\le\frac1{30}$, $k\ge1$, $q>1$,
$v=\operatorname{Sec}(t,1-t,t+x)\le U(t)$ and $N(t,x,k,q,v,\log q)<0$, then $\chi=pq/W>\frac47$.
\end{lemma}

\begin{proof}
Lower $k$ to $1$ (\cref{lem:m2-shape}): $N_1(q)=A_2q^2+A_1q-tp^2W^2-Cq\log q$ with $C=tpBW>0$
and $N_1''=2A_2-C/q$. Suppose $1<q\le q_\chi=4W/(7p)\le\frac87$. The $v$-coefficient of $A_2$ is
$-Y-tY^2<0$, so with $v>\frac34$, $N_1''<-\frac18(7C-16A_2)|_{v=3/4}$, and
$(7C-16A_2)|_{v=3/4}=\sum_{i=0}^4b_i(h)j^i$ with
\[
b_0=-\tfrac{h+1}4(16h^4+23h^3-9h^2-7h+1),\qquad b_1=-\tfrac{h+1}4(87h^3+12h^2-55h-12),
\]
\[
b_2=-\tfrac12(66h^3+33h^2-58h-33),\qquad b_3=\tfrac{67}4(1-h^2),\qquad b_4=\tfrac32(1-h),
\]
all positive for $\frac6{25}\le h\le\frac13$. So $N_1$ is concave on $[1,q_\chi]$. At $q=1$ the secant
bound $v<V_2=1-cp/W^2+(cx/W^2)(\frac12-\frac x3+\frac{x^2}{12}-t)$ gives $N_1(1)>0$ (the cleared value
exceeds a positive multiple of $2h-j^3-4j^2-4j>0$). At $q=q_\chi$, the bounds $V_2$ and $L_+$ give
$N_1(q_\chi)>P(h,j)/(294(15c+11j)p^2)$ with $P=\sum_{i=0}^8P_i(h)j^i$, where $P_0>\frac35$,
$P_1>-13$, $P_2>-26$, $P_3,P_4>0$, $P_5>-330$, $P_6>-585$, $P_7>-140$ and $P_8>0$ on
$\frac6{25}\le h\le\frac13$, each by a one-variable monotonicity or by discarding positive terms; with
$j\le\frac1{30}$ this gives $P>0$. So $N_1>0$ on $[1,q_\chi]$, against $N<0$.
\end{proof}

\begin{lemma}[Casimir slope]\label{lem:m2-slope}
Let $\frac13\le t\le\frac25$, $0<j<\frac14$, $k\ge1$, $q>1$, $\chi\ge\frac47$, $v\ge\frac34$,
$F_{p/W}(q)\le v$, and let $r=qB/W-p$ satisfy $t\le r\le p$. With
$H_d(w)=BL'+\frac epB^2(w-1)+p(1-\chi)/e$ and $L'=\log q+(Wv-c+p^2/W)/p$, one has
$2v-\frac epH_d(w)<0$ for every $w\ge\frac34$; so $\Gamma(w)=(4B^2vw-H_d(w)^2)/(16(k-t^2)^2)$
decreases on $[\frac34,\infty)$.
\end{lemma}

\begin{proof}
$\Gamma'(w)=\frac{B^2}{8(k-t^2)^2}(2v-\frac epH_d(w))$ and $H_d$ increases in $w$, so it suffices to
take $w=\frac34$. With $A=Z/p=1+r/p$ the hypotheses give $\frac{15}{11}\le A<2$ and $A/\chi=B/p^2>2$,
and $j<\frac14$ gives $k-t^2-p^2>\frac7{144}$. The floor bounds $\log q$ below in terms of $v$, and
$2v-\frac epH_d(\frac34)$ is then bounded above by a convex quadratic in one ratio, negative at both
ends of its range
\cite[\texttt{MIXED\_\allowbreak NONSHORT\_\allowbreak FORWARD\_\allowbreak CASIMIR\_\allowbreak MONOTONICITY}]{Notes}.
\end{proof}

\begin{lemma}[Eliminating the event factor]\label{lem:m2-qelim}
Let $\frac13\le t\le\frac25$, $0\le j\le\frac1{30}$, $k\ge1$, $v\le\min\{1,V\}$, $\eta=p/W$,
$\alpha_U=B(w_U-1)/(pW)$ and $\beta_U=pW/B$. Then for $q>1$,
\[
\log q+\alpha_U(q-1)+\beta_U(1/q-1)<F_\eta(q)-U(\eta).
\]
\end{lemma}

\begin{proof}
By \cref{lem:m2-logs}(b) the difference of the two sides exceeds
\[
\frac{y}{6(y+1)(y+2)}(C_0+C_1y+C_2y^2)
\]
with $C_0=12\eta^2-6\eta-12\alpha_U+12\beta_U-12$,
$C_1=18\eta^2-7\eta-18\alpha_U+6\beta_U-12$ and $C_2=6\eta^2-\eta-6\alpha_U-1$. Each $C_i$ increases
in $k$, since $\partial_k\alpha_U<\partial_k\beta_U=-pW/B^2<0$, and decreases in $v$, since
$\partial_v\alpha_U=Y(k/t+Y)/(BpW)>0$; so it suffices to take $k=1$ and $v=V$. There $w_U<1$, so
$\alpha_U<0$, and $\eta\ge\frac12$ gives $C_2=(2\eta-1)(3\eta+1)-6\alpha_U>0$. It remains to show
$4C_0C_2>C_1^2$. With $t=\frac25-A$ and $x=\frac15+2A+j$ on $0\le A\le\frac1{15}$, $0\le j\le\frac1{30}$, the
cleared discriminant is a polynomial in $(A,j)$, concave in $A$ (the positive part of its second
$A$-derivative is dominated using $A^ij^m\le15^{-i}30^{-m}$), and positive on the two edges $j=0$ and
$j=\frac1{30}$, where it is a univariate polynomial whose negative terms are collected against its
positive ones
\cite[\texttt{MIXED\_\allowbreak NONSHORT\_\allowbreak FORWARD\_\allowbreak Q\_\allowbreak ELIMINATION}]{Notes}.
\end{proof}

\begin{lemma}[The envelope Casimir]\label{lem:m2-envelope}
On $\frac13\le t\le\frac25$, $0\le j\le\frac1{30}$, $1\le k\le\frac{31}{30}$, $\frac34\le v\le V$,
\[
\mathcal G=\frac{B^2\{4vw_U-(\phi(1)+v-g)^2\}}{16(k-t^2)^2}>\frac{9193}{50000},
\]
where $\phi(1)=\alpha_U+\beta_U+(Wv-c+p^2/W)/p-p^2/B$.
\end{lemma}

\begin{proof}
The numerator $\mathcal N(v)=B^2\{4vw_U-(\phi(1)+v-g)^2\}$ is a quadratic in $v$. The numerator of
$w_U$ has $v$-derivative $P_v=Y(k/t+Y)$, and $D_1(v)=B(\phi(1)+v-g)$ has $D_1'=B(1+W/p)+P_v/(pW)$,
so the arithmetic--geometric mean inequality and $B>p^2$ give $D_1'^2>4P_v$ and
$\mathcal N''=8P_v-2D_1'^2<0$. Put $t=\frac25-A$, $x=\frac15+2A+j$ and $k=1+Z$ on the box
$0\le A\le\frac1{15}$, $0\le j,Z\le\frac1{30}$. The polynomial $5^{10}\cdot2t^2p^2W^3(-\mathcal N'(\frac34))$
has constant coefficient $6465474$, while the sum of its negative coefficients, each weighted by
the maximum $15^{-i}30^{-m-n}$ of its monomial $A^ij^mZ^n$ on the box, is
$787831672153/787320<6465474$. So $\mathcal N'(\frac34)<0$, $\mathcal N$ decreases on $[\frac34,V]$, and
it remains to show $\mathcal N(V)>16(k-t^2)^2\cdot\frac{9193}{50000}$. The cleared difference is concave
in $A$ by the same domination, and positive on the edges $A=0$ and $A=\frac1{15}$ by collecting
negative coefficients against positive ones
\cite[\texttt{MIXED\_\allowbreak NONSHORT\_\allowbreak FORWARD\_\allowbreak ENVELOPE\_\allowbreak CASIMIR\_\allowbreak SIGN}]{Notes}.
\end{proof}

\runin{The small gap.} By \cref{lem:m2-chi}, the hypotheses of \cref{lem:m2-slope} hold, so
$\gamma^2=\Gamma(w)>\Gamma(w_U)$. At $w_U$ the diagonal divided by $B$ is
$\phi(q)=\log q+\alpha_Uq+\beta_U/q+(Wv-c+p^2/W)/p-p^2/B$, and \cref{lem:m2-qelim} with the floor
gives $0<\phi(q)\le\phi(1)+F_\eta(q)-U(\eta)\le\phi(1)+v-g$; positivity comes from
\cref{lem:m2-slope}, and $U(\eta)=g$. Hence $\Gamma(w_U)\ge\mathcal G>9193/50000$ by
\cref{lem:m2-envelope}, which is the last step of \cref{prop:m2-forward-steps}.
